\scanpage{51}
\margin{Cas transposés l'un de l'autre}
\diagram{\scriptsize
\begin{tabular}{@{}c|c@{}}
\hfill I & \hfill II\\
\begin{tikzcd}[ampersand replacement=\&,column sep=small,row sep=small]
H^1(Y)\arrow[r] \& H^1(X)\\
H^1(A_Y)\arrow[u,"\varphi_Y"']\arrow[r] \& H^1(A_X)\arrow[u,"\varphi_X"']\\
A_Y \& A_X\arrow[l]
\end{tikzcd}
&
\begin{tikzcd}[ampersand replacement=\&,column sep=small,row sep=small]
H^{2m-1}(Y\unsure{.})(m-1)\arrow[r]\arrow[d,"\psi_Y"] \& H^{2n-1}(X)\arrow[d,"\varphi_X"]\\
H^1(B_Y)\arrow[r] \& H^1(B_X)\\
B_Y \& B_X\arrow[l]
\end{tikzcd}\\ \hline
\hfill III & \hfill IV\\
\begin{tikzcd}[ampersand replacement=\&,column sep=small,row sep=small]
H^{2m-1}(Y)(m-1)\arrow[r]\arrow[d,"\psi_Y"] \& H^1(X)\\
H^1(\deleted{\unsure{A}}B_Y)\arrow[r] \& H^1(A_X)\arrow[u,"\varphi_X"']\\
{} \& {}\\[-1ex]
B_Y \& A_X\arrow[l]
\arrow[from=3-1,to=3-2,phantom,"\text{Cas auto-transposé}" description]
\end{tikzcd}
&
\begin{tikzcd}[ampersand replacement=\&,column sep=small,row sep=small]
H^1(Y)\arrow[r,dash,"\unsure{-|\to}" description] \& H^{2n-1}(X)(n-1)\arrow[d,"\psi_X"]\\
H^1(A_Y)\arrow[u,"\varphi_Y"']\arrow[r,dash] \& H^1(B_X)\\
{} \& {}\\[-1ex]
A_Y \& B_X\arrow[l]
\arrow[from=3-1,to=3-2,phantom,"\text{Cas auto-transposé}" description]
\end{tikzcd}\\ \hline
\end{tabular}}

\uline{Cas} I\quad Pour passer des hom de VA à des cycles sur $X\times Y$, il faut inverser $\varphi_Y$

\uline{Cas} II\quad \deleted{\unsure{On peut|Cas pareil}}\hfill --- $\varphi_X$
\margin{$\left.\vphantom{\begin{matrix}1\\1\end{matrix}}\right\}$ il faut des dénominateurs ?}

\uline{Cas} III\quad On peut tjrs passer d'un hom de VA à un cycle : c'est aussi la théorie classique...

\uline{Cas} IV\quad Pour passer des hom de VA à des cycles sur $X\times Y$, il faut inverser $\varphi_Y$ et $\psi_X$, ce qui introduit des dénominateurs.

Pour inverser $\psi_X$, on utilise «th. de Lefschetz». 

Pour inverser $\varphi_X$, on utilise le \unsure{\textquoteleft|\textquotedblleft|\guillemotleft{}}th de Lefschetz'' et le cas particulier des V.A. qui est trivial.
% archive sheet number: 50
